\section{Birational fields and a divisorial cocycle}
\label{pre:section}

We define the field test and the additive birational invariant used in
both factorizations. On a smooth blowup the invariant will be the test
of the center, with a sign determined by the direction of the map.

All varieties and maps are over $\C$, unless a function field is explicitly
specified. A variety is integral. When a finite union is used, every sum is
taken over its integral components. A model of a finitely generated field
$F/\C$ is a variety with function field $F$; a smooth projective model may
always be chosen. A fourfold is \emph{rational} if it is birational over $\C$
to $\mathbb P^4$. Rational connectedness of a field means rational
connectedness of a smooth projective model.
We use characteristic-zero resolution and principalization in the
projective category, originating with Hironaka~\cite{Hironaka64}; precise
forms are given in \cite[Theorems~1.0.1--1.0.3]{Wlodarczyk05}.
Resolving the closure of a joint graph gives a common smooth projective
resolution of finitely many rational maps. Principalization makes the
marked ideals and exceptional boundary into a simple normal-crossing
support; see also \cite[Sections~1.2.2--1.2.3]{AKMW02}.

\subsection{The K3 test}

For a smooth projective variety $V$, write
$V\dashrightarrow R(V)$ for its maximal rationally connected quotient.
Its general fibres are rationally connected, its base is not uniruled,
and the base is determined up to birational equivalence. We use the
standard quotient theorem and the theorem that a fibration with
rationally connected base and general fibre has rationally connected
total space; see \cite[Chapter~IV, \S5]{Kollar96} and
\cite[Corollaries~1.3--1.4]{GHS03}.
The quotient has a proper morphism representative on suitable dense open
subsets \cite[\S0.7 and Theorem~2.3]{Campana92}.

Fix a positive integer $d>2$ and any collection $\mathcal S$ of
isomorphism classes of smooth projective K3 surfaces $S$ such that
\begin{equation}\label{pre:test-collection}
 \operatorname{Pic}(S)=\Z L_S,\qquad L_S\text{ is ample},\qquad
 L_S^2=d,\qquad \operatorname{Aut}_{\C}(S)=\{1\}.
\end{equation}
For odd $d$ the collection is empty. The collection can be restricted
by any further condition, including
an integral polarized transcendental Hodge-isometry condition. Every
degree threshold proved below is uniform in this choice of subcollection.

\begin{definition}[The test function]\label{pre:test}
For a finitely generated field $F/\C$, choose a smooth projective model
$V$ and set
\[
 u(F)=
 \begin{cases}
  1,&\text{if }R(V)\text{ is birational to some }S\in\mathcal S,\\
  0,&\text{otherwise}.
 \end{cases}
\]
For an integral variety $V$, put $u(V)=u(\C(V))$.
\end{definition}

Birational invariance and uniqueness of the quotient make this definition
independent of the model. A K3 surface is not uniruled. Moreover, two
birational smooth projective K3 surfaces are isomorphic. One way to see
the latter assertion is to resolve a birational map as
$S\xleftarrow{p}W\xrightarrow{q}S'$. The pullbacks of the nowhere
vanishing two-forms on $S$ and $S'$ span the same one-dimensional space
$H^0(W,\Omega_W^2)$. Their zero divisors therefore have the same support.
These supports are respectively the exceptional loci of $p$ and $q$.
Every exceptional fibre is connected, so $q$ contracts every fibre of
$p$, and conversely. The image of $(p,q)$ in $S\times S'$ projects
properly, quasi-finitely and birationally to the normal surface $S$,
so that projection is an isomorphism.
Descending in both directions gives inverse morphisms. In particular,
$\operatorname{Bir}_{\C}(S)=\operatorname{Aut}_{\C}(S)$ for a K3 surface.

\begin{lemma}[Rationally connected fibrations]\label{pre:mrc-fibration}
Let $V\dashrightarrow W$ be a dominant rational map whose geometric
generic fibre is integral and has a rationally connected smooth proper
model. Then $R(V)$ and $R(W)$ are birational, and hence $u(V)=u(W)$.
In particular, for any vector bundle $\mathcal E$ of positive rank on
an integral variety $W$,
\[
 u(\mathbb P_W(\mathcal E))=u(W).
\]
The test vanishes on rationally connected varieties and on varieties
of dimension at most one. For an integral surface $W$, it is one
precisely when $W$ is birational to a member of $\mathcal S$.
\end{lemma}

\begin{proof}
Resolve the rational maps and use smooth projective models.
Compose $V\dashrightarrow W$ with the maximal rationally connected
quotient of $W$. A general fibre of the composite fibres over a
rationally connected general fibre of $W\dashrightarrow R(W)$, with
rationally connected general fibre. It is rationally connected by the
fibration theorem of Graber--Harris--Starr
\cite[Corollary~1.3]{GHS03}. Since $R(W)$
is not uniruled, this composite is a maximal rationally connected
quotient of $V$. Indeed, a family of rational curves through general
points of $V$ that is not vertical would induce a covering family of
rational curves on $R(W)$. Uniqueness of the quotient proves the first
assertion. A projective bundle has a projective-space generic fibre.
The remaining statements follow from the dimension of the quotient,
and from the fact that a non-uniruled surface is its own quotient.
\end{proof}

We also use $u$ on finite field extensions and finite Galois sets.
If $k/\C$ is finitely generated and a finite \mbox{\'{e}tale}
$k$-algebra is $A=\prod_i k_i$, define
\begin{equation}\label{pre:finite-orbits}
 u(A)=\sum_i u(k_i).
\end{equation}
Equivalently, for a finite set with continuous action of
$\operatorname{Gal}(\bar k/k)$, sum once over its Galois orbits, using
the finite residue field belonging to each orbit. There is no factor
$[k_i:k]$ in this sum. Each $k_i$ is viewed as a field over $\C$.
Geometric generic fibres over algebraic closures of function fields will
be used to obtain bounds, while $u$ continues to be evaluated on the
original finitely generated fields over $\C$.

\begin{lemma}[Small orbits]\label{pre:small-orbits}
Suppose $k\simeq\C(x,y)$. A finite Galois orbit of size one or two
has test value zero. More generally, if a finite extension $\ell/k$
has a nonidentity $\C$-automorphism, then $u(\ell)=0$.
\end{lemma}

\begin{proof}
The field $k$ is rational. If $u(\ell)=1$, the surface field $\ell$
is the function field of a member of $\mathcal S$, by
Lemma~\ref{pre:mrc-fibration}. A nonidentity field automorphism would
give a nonidentity birational, hence regular, automorphism of that K3
surface. This contradicts the choice of $\mathcal S$. A quadratic
extension in characteristic zero has its nonidentity involution.
\end{proof}

\subsection{Prime divisors as valuations}

Let $X$ be a normal projective variety with function field $F$.
A prime divisor $E\subset X$ determines the normalized discrete
valuation $\operatorname{ord}_E:F^*\to\Z$, whose residue field is
$\C(E)$. Write $\mathcal D_F(X)$ for this set of valuations, using
a specified identification $\C(X)\simeq F$ when necessary. Distinct
valuations are counted separately even when their residue fields are
isomorphic.

For a birational map $f:X\dashrightarrow Y$, identify the function
fields by $f^*: \C(Y)\xrightarrow{\sim}\C(X)=F$ and define
\begin{equation}\label{pre:cocycle-definition}
 c_u(f)=
 \sum_{\nu\in\mathcal D_F(Y)\setminus\mathcal D_F(X)}u(\kappa(\nu))
 -\sum_{\nu\in\mathcal D_F(X)\setminus\mathcal D_F(Y)}u(\kappa(\nu)).
\end{equation}
Thus divisors acquired by the target have positive sign. The target's
valuation set is transported using $f^*$ even when $X=Y$.
This is the divisorial invariant of Lin and Shinder, composed with the
MRC/K3 test \(u\); the sign agrees with their target-acquired minus
source-lost convention
\cite[Definition~2.1, Lemma~2.2 and Proposition~2.3]{LinShinder24}.

\begin{lemma}[Divisorial cocycle]\label{pre:cocycle}
The sums in \eqref{pre:cocycle-definition} are finite. For birational
maps $f:X\dashrightarrow Y$ and $g:Y\dashrightarrow Z$ of normal
projective varieties,
\begin{equation}\label{pre:cocycle-additivity}
 c_u(g\circ f)=c_u(f)+c_u(g),\qquad
 c_u(f^{-1})=-c_u(f).
\end{equation}
An isomorphism in codimension one has value zero.
If $\pi:\widetilde X\to X$ is the blow-up of a smooth projective
variety along a smooth center with connected components $Z_i$ of
codimension at least two, then
\begin{equation}\label{pre:blowup-sign}
 c_u(\pi^{-1})=\sum_i u(Z_i),\qquad
 c_u(\pi)=-\sum_i u(Z_i).
\end{equation}
\end{lemma}

\begin{proof}
Choose dense open subsets on which $f$ is an isomorphism. Every divisor
whose generic point belongs to these opens has a corresponding divisor
on the other model with exactly the same valuation. Any valuation in
the symmetric difference must therefore be represented by a divisorial
component of one of the two closed complements. There are finitely many
such components. This also proves that an isomorphism in codimension
one has value zero.

For the composition use the common field $F=\C(X)$, transporting
$\C(Y)$ by $f^*$ and $\C(Z)$ by $f^*g^*$. Put
$D_X=\mathcal D_F(X)$, $D_Y=\mathcal D_F(Y)$ and
$D_Z=\mathcal D_F(Z)$. Their pairwise symmetric differences are finite.
On the set of all prime divisorial valuations of $F$,
\[
 \mathbf 1_{D_Z}-\mathbf 1_{D_X}
 = (\mathbf 1_{D_Y}-\mathbf 1_{D_X})
   +(\mathbf 1_{D_Z}-\mathbf 1_{D_Y}).
\]
All three differences have finite support. Multiplying by
$u(\kappa(\nu))$ and summing proves additivity. Transport by a
$\C$-field isomorphism identifies residue fields, so the last term
is precisely $c_u(g)$. Reversing the two valuation sets gives the
inverse formula.

For the blow-up, the target of the extraction
$\pi^{-1}:X\dashrightarrow\widetilde X$ retains every original
divisorial valuation and adds exactly one exceptional divisor
$E_i=\mathbb P_{Z_i}(N_{Z_i/X})$ for each $Z_i$. Its residue field is
purely transcendental over $\C(Z_i)$, of transcendence degree
$\operatorname{codim}_X(Z_i)-1$. Lemma~\ref{pre:mrc-fibration} gives
$u(E_i)=u(Z_i)$. The stated signs now follow from the definition and
the inverse formula.
\end{proof}

For the generic-surface interpretation, let $V\to Q$ have a smooth
surface generic fibre over $k=\C(Q)$. Blowing up a closed point $P$ of
that fibre creates the exceptional curve $\mathbb P^1_{\kappa(P)}$.
Its function field is $\kappa(P)(t)$, so the corresponding divisorial
valuation of $\C(V)$ contributes $u(\kappa(P))$ once, by
Lemma~\ref{pre:mrc-fibration}. There is no factor $[\kappa(P):k]$.
This is the compatibility between \eqref{pre:finite-orbits} and the
divisorial count used for generic surface links.

\subsection{The smooth factorization used below}

\begin{theorem}[Weak factorization]\label{pre:weak-factorization}
A birational map between smooth projective varieties over $\C$ factors
into a finite zigzag of blow-ups and blow-downs along smooth connected
centers, with all intermediate varieties smooth and projective.
Nontrivial steps may be taken to have centers of codimension at least two.
\end{theorem}

This is the projective case of
\cite[Theorem~0.1.1]{AKMW02}; blowing up a smooth divisor is an
isomorphism and may be omitted. In a factorization of a birational map
from $\mathbb P^4$, every intermediate variety is itself a smooth
projective rational fourfold, hence rationally connected. Its centers
have dimension zero, one, or two. Lemma~\ref{pre:cocycle} computes
$c_u$ as the signed sum of their tests along the zigzag; only surface
centers can have nonzero test.
